\section{Background and Related Work}
\label{sec:related}

Deciding that one candidate patch needs less validation work than the next is
not new, and \cref{tab:positioning} sorts the ways it has been done by what the
decision is computed from and what that demands of the candidate space. We put
\CEGMem\ in the same row as a mechanism a decade older than LLM repair, and
then say what is left over.

\begin{table*}[t]
\centering
\scriptsize
\caption{Where \CEGMem\ sits. The columns are factual properties of each
mechanism, not a scoring of it. \CEGMem\ shares the last column's value with a
line a decade older than LLM repair; what is new is not a cheaper
validation---a classical prioritizer beats ours on executions
(\cref{sec:results-policies})---but that the same evidence removes oracle
invocations, which no row above does, from a key that needs nothing of the
candidate space.}
\label{tab:positioning}
\setlength{\tabcolsep}{4pt}
\renewcommand{\arraystretch}{1.05}
\begin{tabular}{@{}L{2.1cm}L{4.0cm}L{4.6cm}L{3.4cm}L{2.5cm}@{}}
\toprule
\textbf{Line of work} & \textbf{Representatives} &
\textbf{The decision is computed from} &
\textbf{Requires of the candidate space} &
\textbf{Effect on validation work} \\
\midrule
Candidate equivalence &
AE~\cite{weimer2013leveraging}; f1x~\cite{mechtaev2018testequivalence} &
candidates related to \emph{each other}: syntactically, or by the values a
modified expression takes &
enumerable; an edit shape the analysis models &
removes \\

Validation engineering &
UniAPR~\cite{chen2021uniapr}; ExpressAPR~\cite{xiao2024expressapr} &
a warm JVM; online patch equivalence, schemata, virtualization &
mutation-shaped patches, for the equivalence part &
accelerates; ExpressAPR also removes \\

Test prioritization &
fault-recorded~\cite{qi2013efficient}; modification-point
aware~\cite{venugopal2020modification}; AE's
TestStrat~\cite{weimer2013leveraging}; RTS study~\cite{lou2024when} &
which of the project's own tests killed earlier candidates, keyed on nothing or
on the edit location &
nothing &
reorders, stops early \\

LLM repair agents &
RepairAgent~\cite{bouzenia2025repairagent}; SWE-agent~\cite{yang2024sweagent};
Agentless~\cite{xia2025agentless}; ChatRepair~\cite{xia2024chatrepair};
ExpeRepair~\cite{mu2026experepair}; \cite{vo2026merit} &
a flat transcript, or demonstrations and insights, retrieved into the prompt &
nothing &
none \\

Counterexample-guided repair &
\Sketch~\cite{solarlezama2006sketching}; LLM
CEGIS~\cite{orvalho2025counterexample} &
counterexamples, as prompt content for the next attempt &
a complete synthesizer (\Sketch); nothing (LLM) &
none \\
\midrule
\textbf{\CEGMem} & this paper &
counterexamples \emph{retained across candidates} and re-executed, indexed by
edit location $\loc$ &
nothing &
reorders, stops early; \emph{removes oracle invocations} \\
\bottomrule
\end{tabular}
\end{table*}


\subsection{Removing validations: candidate equivalence}
\label{sec:related-avoid}
One family relates candidates \emph{to each other} and validates one
representative per class. AE quotients the patch space under an approximate
program-equivalence relation: on the $45$ defects it and GenProg both repair,
from a $105$-defect benchmark, it needs $186$ test-suite evaluations against
$3{,}252$~\cite{weimer2013leveraging}. Test-equivalence analysis partitions
candidates by the values a modified expression evaluates to during a test, so
one execution decides an entire class~\cite{mechtaev2018testequivalence}, and
ExpressAPR detects that equivalence \emph{online}, reporting
${\approx}137\times$ over plain validation on
Defects4J~\cite{xiao2024expressapr}. UniAPR, by contrast, purely
accelerates---it redefines bytecode inside a warm JVM and still runs every
test~\cite{chen2021uniapr}. The equivalence mechanisms give a stronger
guarantee than ours, and they need what an LLM does not supply: an enumerable
candidate space and an edit shape the analysis can model. That is why they do
not transfer, and it is the only reason.

\subsection{Reordering validations, and where our guard belongs}
\label{sec:related-order}
The other family validates every candidate and changes the order in which
evidence against it is sought. Fault-recorded prioritization runs the tests
that killed earlier candidates first~\cite{qi2013efficient};
modification-point--aware prioritization keeps one ordered suite per edit
location~\cite{venugopal2020modification}; AE's TestStrat---in the same paper as
its quotient space---takes the test likeliest to fail and stops
there~\cite{weimer2013leveraging}; and a study of \num{2532915} patches from
$12$ repair systems on $395$ Defects4J bugs quantifies what selection and
prioritization save~\cite{lou2024when}.

\textbf{Our guard belongs here, and we state the identity rather than argue
with it.} Inline the guard into the oracle---stored counterexamples first, stop
at the first failure---and the executions performed are the same ones, in the
same order, at the same cost. Nothing is \emph{removed}; a validation is made
to terminate on its first case. We therefore claim an ordering effect and
report it as one: in \armtyped\ a round the guard resolves costs $1.13$
program executions where a round reaching the oracle costs $18.9$, and memory
turns roughly $7.6$ oracle calls per episode into roughly $7.4$
early-terminating ones (\cref{sec:results-cost}). What is \emph{not} shared with prior work is the key
and the evidence. Those tables rank the project's own tests by a failure
history over candidates; ours holds inputs discovered inside the episode by
candidates the oracle already refuted, re-executed against a candidate nobody
has run, indexed by a location a diff supplies. The per-modification-point
table of~\cite{venugopal2020modification} is the nearest neighbour, and we run
it: over byte-identical candidates it reaches a refuting case in fewer
executions than our guard on $80$ of $99$ tasks, by a median $4.6\%$
(\cref{sec:results-policies}). What separates the two is not the cost of a
validation but whether one happens: no prioritizer in this line moves oracle
invocations.

\subsection{Counterexample-guided repair, and why replay is not vacuous}
\label{sec:related-cegis}
That a counterexample should constrain the next candidate is the engine of the
loop \Sketch\ introduced~\cite{solarlezama2006sketching}, and the loop has been
run with an LLM in the synthesizer's place over $1{,}431$ student C
programs~\cite{orvalho2025counterexample}. There, retaining counterexamples and
replaying them would be \emph{vacuous}: a constrained synthesizer cannot emit a
candidate one of them refutes, so a guard would never fire. Losing that
constraint is what creates the opportunity, and an LLM takes it often. In
\armguard, $24.2\%$ of proposals are byte-identical to an earlier one in the
same episode and $30.2\%$ of the rounds the guard resolves are such a repeat;
of those resolutions $12.8\%$ re-propose a patch the store already holds, and
the remaining $87.2\%$ are patches the store never held, refuted by a stored
input anyway (\appref{tab:related}{the supplement}). Neither case can arise in
\Sketch. Orvalho et al.\ use the counterexample as
prompt content for the next attempt---our \armsteer\ arm, the half we find
null; we add the second use of the same evidence.

\ifsuppmatter\else\begin{table}[t]
\centering
\scriptsize
\caption{What makes replaying a stored counterexample non-vacuous, and what the
$\loc$ index does not buy. Every quantity is derived by replaying the frozen
round log---no program was executed to produce this table
(\texttt{tools/related\_addenda.py}). The upper block is the opportunity an LLM
proposer creates and a constrained synthesizer does not; the lower block is why
we claim no benefit for indexing on this corpus. The last two rows of the upper
block partition the resolutions; the repeat-proposal row above them cuts across
that partition.}
\label{tab:related}
\setlength{\tabcolsep}{4pt}
\renewcommand{\arraystretch}{1.06}
\begin{tabular}{@{}L{4.15cm}rrr@{}}
\toprule
& \armuntyped & \armtyped & \armguard \\
\midrule
\multicolumn{4}{@{}l}{\textit{The proposer repeats itself}} \\
Rounds the guard resolved & 3785 & 3649 & 2253 \\
Proposals made by the arm & 4732 & 4557 & 2807 \\
\quad byte-identical to an earlier one & 24.3\% & 27.1\% & 24.2\% \\
Of resolutions, a repeat proposal & 30.4\% & 33.9\% & 30.2\% \\
\quad a patch the store already held & 12.6\% & 13.8\% & 12.8\% \\
\quad a patch the store never held & 87.4\% & 86.2\% & 87.2\% \\
\midrule
\multicolumn{4}{@{}l}{\textit{Memory is too small for order to matter}} \\
Entries held at resolution, median & 1 & 1 & 1 \\
\quad maximum observed & 5 & 4 & 3 \\
Resolutions with one entry or none & 57.5\% & 63.3\% & 60.5\% \\
Resolutions after one consultation & 83.3\% & 87.6\% & 86.4\% \\
\midrule
\multicolumn{4}{@{}l}{\textit{Some evidence is found, not given}} \\
Distinct stored inputs per episode & 1.53 & 1.44 & 1.47 \\
Stored inputs beyond the round-1 one & 34.7\% & 30.5\% & 32.0\% \\
\bottomrule
\end{tabular}
\end{table}
\fi

\subsection{Memory in LLM repair agents}
LLMs repair code by conversing: ChatRepair iterates on test
failures~\cite{xia2024chatrepair}, agentic systems localize, edit and validate
over many turns~\cite{bouzenia2025repairagent,zhang2024autocoderover,%
yang2024sweagent}, ``agentless'' pipelines keep a fixed propose--validate
workflow~\cite{xia2025agentless}, and general frameworks add verbal
self-reflection~\cite{shinn2023reflexion,chen2024selfdebug}. In all of them the
record of past attempts is a flat transcript and the cost model is dollars or
tokens. The closest memory-augmented system,
ExpeRepair~\cite{mu2026experepair}, retrieves demonstrations and reflective
insights to condition generation but never validation; nothing in that index
recognizes that a new candidate is refuted by an input already in hand. Vo et
al.~\cite{vo2026merit} type failures for text-to-SQL repair, conditioning
\emph{retrieval for the prompt} where we condition \emph{the order of
validation}, and report their typed variant ``is not reliably separated from
untyped dynamic retrieval''---the phenomenon we measure directly. The guard is
a property of memory, not of a controller, so it is orthogonal to every agent
above. Finally, since a patch can pass the visible tests and still be
wrong~\cite{xiong2018patchsim} and
LLM-era benchmarks must worry about leakage, we use
ConDefects~\cite{wu2024condefects} rather than
Defects4J~\cite{just2014defects4j} or SWE-bench~\cite{jimenez2024swebench}, and
treat acceptance as a claim about $\Gk$ audited against $\Gpool$
(\cref{sec:results-depth}).
